\chapter{Boolean Analysis — Circuit Helpers}
%%% generated by the blueprint pipeline from TCSlib.BooleanAnalysis.LMN.CircuitHelpers
\section{Overview}

This module collects helper lemmas and constructions used to prove switching-lemma
bounds for depth-$2$ Boolean circuits in the LMN development. It provides a cleaning
procedure that removes contradictory and duplicated literals from DNF/CNF formulas
(while preserving their evaluation and not increasing width), general switching-lemma
bounds that drop the syntactic ``no duplicate variable'' requirement, and translations
from shallow Boolean circuits into DNF/CNF formulas.

%%%%%%%%%%%%%%%%%%%%%%%%%%%%%%%%%%%%%%%%%%%%%%%%%%%%%%%%%%%%%%%%%%%%%%%%%%%%%%%
\section{Declarations}

\begin{lemma}[Restriction respects pointwise equality]
\lean{LMN.restrictFn_ext'}
\label{LMN.restrictFn_ext'}
\leanok
\uses{SwitchingLemma2.Restriction, SwitchingLemma2.restrictFn}
\difficulty{1}
Let $f, g : (\mathrm{Fin}\,n \to \mathrm{Bool}) \to \mathrm{Bool}$ be two Boolean
functions that agree at every point, so that $f(x) = g(x)$ for all $x$. Then for every
restriction $\rho$ on $n$ variables, the restricted functions coincide: the restriction
of $f$ by $\rho$ equals the restriction of $g$ by $\rho$.
\end{lemma}

\begin{lemma}[Depth-exceedance probability depends only on the function]
\lean{LMN.bernoulliRestrProb_congr_fn}
\label{LMN.bernoulliRestrProb_congr_fn}
\leanok
\uses{LMN.dedupTermVar, LMN.restrictFn_ext', SwitchingLemma2.bernoulliRestrProb, SwitchingLemma2.restrictFn, dtDepth}
\difficulty{3}
Let $f, g : (\mathrm{Fin}\,n \to \mathrm{Bool}) \to \mathrm{Bool}$ be Boolean functions
on $n$ variables that agree at every input, and fix a real parameter $p$ and a threshold
$t \in \bbn$. Then the Bernoulli($p$) restriction probability of the event that a random
restriction $\rho$ produces a restricted function of decision-tree depth strictly
greater than $t$ is the same for $f$ as for $g$.
\end{lemma}

\begin{definition}[De-duplicate literals by variable]
\lean{LMN.dedupTermVar}
\label{LMN.dedupTermVar}
\leanok
\uses{Term}
Given a term $t$ (a list of literals), $\mathrm{dedupTermVar}\,t$ keeps each literal only
if no earlier-retained literal already uses the same variable, removing literals whose
variable already appeared.
\end{definition}

\begin{definition}[Contradictory term test]
\lean{LMN.termHasContradiction}
\label{LMN.termHasContradiction}
\leanok
\uses{Term}
$\mathrm{termHasContradiction}\,t$ returns \texttt{true} iff the term $t$ contains two
literals on the same variable with opposite polarities (one negated, one not).
\end{definition}

\begin{definition}[Clean a DNF]
\lean{LMN.cleanDNF}
\label{LMN.cleanDNF}
\leanok
\uses{DNF, LMN.dedupTermVar, LMN.termHasContradiction}
$\mathrm{cleanDNF}\,d$ first discards every contradictory term of the DNF $d$ and then
applies $\mathrm{dedupTermVar}$ to each remaining term.
\end{definition}

\begin{definition}[Clean a CNF]
\lean{LMN.cleanCNF}
\label{LMN.cleanCNF}
\leanok
\uses{CNF, LMN.dedupTermVar, LMN.termHasContradiction}
$\mathrm{cleanCNF}\,c$ filters out the contradictory (tautological) clauses of the CNF $c$
and then de-duplicates the literals of each remaining clause via $\mathrm{dedupTermVar}$.
\end{definition}

\begin{lemma}[De-duplication produces a repetition-free term]
\lean{LMN.dedupTermVar_nodup}
\label{LMN.dedupTermVar_nodup}
\leanok
\uses{LMN.dedupTermVar, Literal, Term}
\difficulty{4}
Let $t$ be a term over $n$ variables, that is, a list of literals, each a variable in
$\{0,\dots,n-1\}$ together with a negation flag. Form a new term $t'$ by scanning $t$
and retaining a literal exactly when no already-retained literal has the same variable,
discarding any literal whose variable has already occurred. Then $t'$ has no repeated
literal: no literal occurs twice in $t'$.
\end{lemma}

\begin{lemma}[Variable-injectivity of literal de-duplication]
\lean{LMN.dedupTermVar_var_inj}
\label{LMN.dedupTermVar_var_inj}
\leanok
\uses{LMN.dedupTermVar, Literal, Term}
\difficulty{4}
Let $t$ be a term over $n$ variables, that is, a list of literals, each literal
consisting of a variable in $\{0,1,\dots,n-1\}$ together with a negation flag. Let $t'$
be the de-duplicated term obtained by scanning $t$ and retaining each literal only when
no earlier-retained literal already uses the same variable. Then within $t'$ the
variable of a literal determines that literal: any two literals $l_1, l_2 \in t'$ whose
variables agree are equal, that is, they have the same variable and the same negation
flag.
\end{lemma}

\begin{lemma}[De-duplication does not increase term length]
\lean{LMN.dedupTermVar_width_le}
\label{LMN.dedupTermVar_width_le}
\leanok
\uses{LMN.contradiction_term_eval_false, LMN.dedupTermVar, Term}
\difficulty{4}
Let $t$ be a term, that is, a finite list of literals over the variables $\{0, 1, \dots,
n-1\}$, and let $t'$ be the term obtained from $t$ by scanning it and keeping each
literal only when no literal already kept uses the same variable. Then $t'$ has length
at most that of $t$.
\end{lemma}

\begin{lemma}[Contradictory terms evaluate to false]
\lean{LMN.contradiction_term_eval_false}
\label{LMN.contradiction_term_eval_false}
\leanok
\uses{LMN.termHasContradiction, Literal, Literal.eval, Term, Term.eval}
\difficulty{4}
Let $t$ be a term (a conjunction of literals) over the variables $x_0,\dots,x_{n-1}$. If
$t$ contains two literals on the same variable with opposite polarities — one negated
and one not — then $t$ evaluates to false under every assignment $x : \{0,\dots,n-1\}
\to \{\mathrm{true},\mathrm{false}\}$.
\end{lemma}

\begin{lemma}[De-duplication preserves term value on consistent terms]
\lean{LMN.dedupTermVar_preserves_term_eval}
\label{LMN.dedupTermVar_preserves_term_eval}
\leanok
\uses{LMN.dedupTermVar, LMN.termHasContradiction, Literal, Literal.eval, Term, Term.eval}
\difficulty{5}
Let $t$ be a term over $n$ variables, that is, a list of literals $l$, each carrying a
variable $l.\mathrm{var} \in \{0,\dots,n-1\}$ and a polarity flag, and interpreted
conjunctively: $t$ evaluates to true on an assignment $x \in \{0,1\}^n$ exactly when
every literal in $t$ holds under $x$. Suppose $t$ is consistent, meaning it contains no
two literals on the same variable with opposite polarities. Let $t'$ be the term
obtained from $t$ by keeping, for each variable, only its first occurrence and
discarding every later literal whose variable has already appeared. Then $t'$ and $t$
have the same value on every assignment $x$.
\end{lemma}

\begin{lemma}[Cleaning a DNF preserves its Boolean function]
\lean{LMN.cleanDNF_eval}
\label{LMN.cleanDNF_eval}
\leanok
\uses{DNF, DNF.eval, LMN.cleanDNF, LMN.contradiction_term_eval_false, LMN.dedupTermVar_preserves_term_eval, LMN.termHasContradiction}
\difficulty{4}
Let $d$ be a DNF formula in $n$ Boolean variables — a disjunction of terms, each term a
conjunction of literals over the variables $\{0,\dots,n-1\}$ — and form the cleaned
formula by first discarding every term that contains some variable together with its
negation, and then, in each surviving term, keeping only the first literal on each
variable and deleting any later literal repeating that variable. Then the cleaned
formula and $d$ compute the same Boolean function: for every assignment $x \colon
\{0,\dots,n-1\} \to \{\mathrm{true},\mathrm{false}\}$, the cleaned formula evaluates to
true at $x$ if and only if $d$ does.
\end{lemma}

\begin{lemma}[Contradictory clauses are tautologies]
\lean{LMN.contradiction_clause_eval_true}
\label{LMN.contradiction_clause_eval_true}
\leanok
\uses{CNF.evalClause, LMN.termHasContradiction, Literal.eval, Term}
\difficulty{2}
Let $t$ be a term over $n$ Boolean variables — a list of literals, each a variable
together with a polarity — and suppose $t$ contains two literals on the same variable
with opposite polarities, one negated and one not. Then, read as a single clause under
disjunction semantics, $t$ evaluates to true under every assignment $x : \{0,\dots,n-1\}
\to \{\mathrm{true},\mathrm{false}\}$; that is, some literal of $t$ holds under $x$, for
every $x$.
\end{lemma}

\begin{lemma}[De-duplication preserves clause value]
\lean{LMN.dedupTermVar_preserves_clause_eval}
\label{LMN.dedupTermVar_preserves_clause_eval}
\leanok
\uses{CNF.evalClause, LMN.cleanDNF_width_le, LMN.dedupTermVar, LMN.termHasContradiction, Literal, Literal.eval, Term}
\difficulty{5}
Let $t$ be a term over the variables $x_0,\dots,x_{n-1}$, that is, a list of literals,
and let $x \in \{0,1\}^n$ be an assignment. Suppose $t$ contains no two literals on the
same variable with opposite polarities. Let $t'$ be the term obtained from $t$ by
scanning it and keeping a literal only when no earlier-kept literal uses the same
variable, so that each variable of $t$ occurs at most once in $t'$. Then $t'$ and $t$
have the same value as a clause under $x$: some literal of $t'$ is satisfied by $x$ if
and only if some literal of $t$ is.
\end{lemma}

\begin{lemma}[Cleaning a CNF preserves its truth value]
\lean{LMN.cleanCNF_eval}
\label{LMN.cleanCNF_eval}
\leanok
\uses{CNF, CNF.eval, LMN.cleanCNF, LMN.contradiction_clause_eval_true, LMN.dedupTermVar_preserves_clause_eval, LMN.termHasContradiction}
\difficulty{4}
Represent a CNF formula over the variables indexed by $\{0,1,\dots,n-1\}$ as a
conjunction of clauses, each clause being a disjunction of literals, where a literal is
a variable or its negation. Cleaning such a formula proceeds in two steps: first delete
every clause that contains some variable together with both polarities (one literal
negated, one not), and then, within each surviving clause, keep a literal only when its
variable has not already occurred among the literals retained earlier in that clause.
Then for every CNF formula $c$ and every Boolean assignment $x \colon
\{0,\dots,n-1\}\to\{\mathrm{true},\mathrm{false}\}$, the cleaned formula and $c$ take
the same Boolean value at $x$.
\end{lemma}

\begin{lemma}[Cleaning a DNF does not increase its width]
\lean{LMN.cleanDNF_width_le}
\label{LMN.cleanDNF_width_le}
\leanok
\uses{DNF, DNF.width, LMN.cleanDNF, LMN.dedupTermVar, LMN.dedupTermVar_width_le, LMN.termHasContradiction, Term.width}
\difficulty{5}
Let $d$ be a DNF formula over $n$ variables — a disjunction of terms, each term a
conjunction of literals — and form its cleaned version by first deleting every
contradictory term (one containing a variable together with its negation) and then, in
each surviving term, discarding any literal whose variable already occurs in an earlier
retained literal of that term. The width of a term is its number of literals, and the
width of a formula is the largest width of any of its terms (with width $0$ for the
empty formula). Then the width of the cleaned formula is at most the width of $d$.
\end{lemma}

\begin{lemma}[Cleaning a CNF does not increase its width]
\lean{LMN.cleanCNF_width_le}
\label{LMN.cleanCNF_width_le}
\leanok
\uses{CNF, CNF.width, LMN.cleanCNF, LMN.dedupTermVar, LMN.dedupTermVar_width_le, LMN.termHasContradiction, Term, Term.width}
\difficulty{5}
Let $c$ be a CNF formula, that is, a list of clauses, each clause being a list of
literals over the variables $\{0,\dots,n-1\}$. Define the width of a clause to be its
number of literals and the width of $c$ to be the maximum width among its clauses (and
$0$ if $c$ has no clauses). Form the cleaned formula from $c$ by discarding every clause
that contains some variable together with both of its polarities, and then, within each
surviving clause, keeping only the first occurrence of each variable and deleting all
later literals on a variable already seen. Then the width of the cleaned formula is at
most the width of $c$.
\end{lemma}

\begin{lemma}[Variable-injectivity of literals in a cleaned DNF]
\lean{LMN.cleanDNF_var_inj}
\label{LMN.cleanDNF_var_inj}
\leanok
\uses{DNF, LMN.cleanDNF, LMN.dedupTermVar, LMN.dedupTermVar_var_inj}
\difficulty{3}
Let $d$ be a DNF formula over $n$ variables — a list of terms, each term a list of
literals, where a literal names a variable in $\{0,\dots,n-1\}$ together with a polarity
(negated or not). Form a new formula $d'$ from $d$ by first discarding every term that
contains two literals on the same variable with opposite polarities, and then, in each
surviving term, keeping only the first literal encountered on each variable and deleting
all later literals that repeat an already-kept variable. Then in every term of $d'$ the
variable determines the literal: for any two literals $l_1, l_2$ occurring in such a
term, if $l_1$ and $l_2$ name the same variable, then $l_1 = l_2$.
\end{lemma}

\begin{lemma}[Cleaned DNF terms have no repeated literals]
\lean{LMN.cleanDNF_nodup}
\label{LMN.cleanDNF_nodup}
\leanok
\uses{DNF, LMN.cleanDNF, LMN.dedupTermVar_nodup}
\difficulty{2}
Let $d$ be a DNF formula in $n$ variables, and clean it by first deleting every term
that contains some variable together with its negation, and then, within each surviving
term, keeping only the first literal on each variable and discarding any later literal
whose variable has already appeared. Every term of the resulting formula is then
duplicate-free: no literal occurs twice in it, where two literals count as equal when
they share both variable and polarity.
\end{lemma}

\begin{lemma}[Variable-injectivity of cleaned clauses]
\lean{LMN.cleanCNF_var_inj}
\label{LMN.cleanCNF_var_inj}
\leanok
\uses{CNF, LMN.cleanCNF, LMN.dedupTermVar_var_inj}
\difficulty{2}
Let $c$ be a CNF formula over the variable set $\{0,\dots,n-1\}$, and let $\hat c$ be
the formula obtained from $c$ by first deleting every clause that contains some variable
together with its negation, and then, in each surviving clause, keeping only one literal
per variable. Then within any single clause of $\hat c$ the variable of a literal
determines that literal completely: if two literals of the clause share the same
variable, they are equal (and in particular have the same polarity).
\end{lemma}

\begin{lemma}[Cleaned clauses have no repeated literals]
\lean{LMN.cleanCNF_nodup}
\label{LMN.cleanCNF_nodup}
\leanok
\uses{CNF, LMN.cleanCNF, LMN.dedupTermVar_nodup}
\difficulty{2}
Let $c$ be a CNF formula over $n$ variables, and form its cleaned version by first
discarding every clause that contains some variable both positively and negatively, and
then deleting from each surviving clause every literal whose variable already occurs
earlier in that clause. Then each clause of the resulting formula, viewed as a list of
literals, contains no repeated literal.
\end{lemma}

\begin{theorem}[General switching lemma for DNFs]
\lean{LMN.switching_bernoulli_dtDepth_dnf_general}
\statementsource{boolean-ch04-dnf-switching-lmn}{pdf-fc447f326910-p008-b005}
\statementsource{switching-lemma}{
  pdf-b5e074215b9e-p001-b008,
  pdf-b5e074215b9e-p002-b001,
  pdf-b5e074215b9e-p002-b002,
  pdf-b5e074215b9e-p002-b003,
  pdf-b5e074215b9e-p002-b004
}
\label{LMN.switching_bernoulli_dtDepth_dnf_general}
\leanok
\uses{DNF, DNF.eval, DNF.width, LMN.cleanDNF, LMN.cleanDNF_eval, LMN.cleanDNF_nodup, LMN.cleanDNF_var_inj, LMN.cleanDNF_width_le, SwitchingBernoulli.switching_bernoulli_dtDepth_dnf, SwitchingLemma2.bernoulliRestrProb, SwitchingLemma2.restrictFn, dtDepth}
\difficulty{2}
Let $f$ be a DNF formula on $n \ge 1$ Boolean variables whose width is at most $w$,
where $w \ge 1$, and let $p \in \bbr$ satisfy $0 < p \le 1$ together with $p \le
\tfrac{1}{40w}$. Draw a Bernoulli($p$) random restriction $\rho$, which independently
leaves each of the $n$ variables free with probability $p$ and otherwise fixes it to a
uniformly random Boolean value, and let $f_\rho$ denote the function of the free
variables obtained by evaluating $f$ with the fixed coordinates set according to $\rho$.
Then for every $t \in \bbn$,
\[
  \bbP_\rho\!\left[\operatorname{depth}(f_\rho) > t\right]
  \;\le\; \left(\tfrac12\right)^{t} + \exp\!\left(-\tfrac{np}{3}\right),
\]
where $\operatorname{depth}(f_\rho)$ is the decision-tree depth of $f_\rho$, i.e.\ the
least depth of a decision tree computing it.
\end{theorem}

\begin{theorem}[General switching lemma for CNF formulas]
\lean{LMN.switching_bernoulli_dtDepth_cnf_general}
\statementsource{boolean-ch04-dnf-switching-lmn}{pdf-fc447f326910-p008-b005}
\statementsource{switching-lemma}{
  pdf-b5e074215b9e-p001-b008,
  pdf-b5e074215b9e-p003-b001
}
\label{LMN.switching_bernoulli_dtDepth_cnf_general}
\leanok
\uses{CNF, CNF.eval, CNF.width, LMN.cleanCNF, LMN.cleanCNF_eval, LMN.cleanCNF_nodup, LMN.cleanCNF_var_inj, LMN.cleanCNF_width_le, SwitchingBernoulli.switching_bernoulli_dtDepth_cnf, SwitchingLemma2.bernoulliRestrProb, SwitchingLemma2.restrictFn, dtDepth}
\difficulty{2}
Let $f$ be a CNF formula in $n \ge 1$ Boolean variables whose width is at most $w$, with
$w \ge 1$, and let $p \in \bbr$ satisfy $0 < p \le 1/(40w)$ and $p \le 1$. Consider a
Bernoulli($p$) random restriction $\rho$ of the $n$ variables, in which each variable is
independently left free with probability $p$ and otherwise fixed to $\mathtt{false}$ or
to $\mathtt{true}$, each with probability $(1-p)/2$, and write $f_\rho$ for the Boolean
function on the free coordinates obtained by evaluating $f$ with the fixed coordinates
set by $\rho$. Then for every $t \in \bbn$,
\[
\Pr_{\rho}\bigl[\,\mathrm{depth}(f_\rho) > t\,\bigr] \;\le\; \left(\tfrac12\right)^{t} +
\exp\!\left(-\,\tfrac{np}{3}\right),
\]
where $\mathrm{depth}(f_\rho)$ denotes the decision-tree depth of $f_\rho$.
\end{theorem}

\begin{lemma}[Depth-one gates have only literal children]
\lean{LMN.depth_le_one_children_are_lits}
\label{LMN.depth_le_one_children_are_lits}
\leanok
\uses{BoolCircuit.Circuit, BoolCircuit.Circuit.depth, BoolCircuit.Lit}
\difficulty{4}
Let $C$ be a Boolean circuit on $n$ variables of the form $g(c_1,\dots,c_k)$, a single
gate $g$ (either AND or OR) applied to a list of child circuits $c_1,\dots,c_k$. If the
depth of $C$ is at most $1$, then every child $c_i$ is a literal, that is, $c_i =
\ell_i$ for some literal $\ell_i$ (a variable index together with a polarity).
\end{lemma}

\begin{lemma}[Children of a depth-two node have depth at most one]
\lean{LMN.depth_le_two_children_depth_le_one}
\label{LMN.depth_le_two_children_depth_le_one}
\leanok
\uses{BoolCircuit.Circuit, BoolCircuit.Circuit.depth}
\difficulty{2}
Let $g = \mathrm{node}(\mathit{isAnd}, cs)$ be a Boolean circuit consisting of a single
gate (an AND gate if $\mathit{isAnd}$ is true, an OR gate otherwise) applied to a list
$cs$ of child circuits, and suppose its depth satisfies $\mathrm{depth}(g) \le 2$. Then
every child circuit $c$ in $cs$ has depth $\mathrm{depth}(c) \le 1$.
\end{lemma}

\begin{definition}[Depth-$\le 1$ AND subcircuit to a term]
\lean{LMN.depth1AndToTerm}
\label{LMN.depth1AndToTerm}
\leanok
\uses{BoolCircuit.Circuit, BoolCircuit.Lit.toLiteral, Term}
Converts a depth-$\le 1$ AND subcircuit to a term (a conjunction): a single literal becomes
the one-element term, and an AND node becomes the list of its literal children translated
via $\mathrm{toLiteral}$.
\end{definition}

\begin{definition}[Depth-$\le 2$ OR-top circuit to a DNF]
\lean{LMN.depth2OrToDNF}
\label{LMN.depth2OrToDNF}
\leanok
\uses{BoolCircuit.Circuit, BoolCircuit.Lit.toLiteral, DNF}
Converts a list of children $cs$ of a top OR gate (of depth $\le 2$) into a DNF: each literal
child becomes a singleton term, each AND child becomes the term of its literals, and each OR
child contributes the singleton terms of its literals.
\end{definition}

\begin{definition}[Depth-$\le 2$ AND-top circuit to a CNF]
\lean{LMN.depth2AndToCNF}
\label{LMN.depth2AndToCNF}
\leanok
\uses{BoolCircuit.Circuit, BoolCircuit.Lit.toLiteral, CNF}
Converts a list of children $cs$ of a top AND gate (of depth $\le 2$) into a CNF: each literal
child becomes a singleton clause, each OR child becomes the clause of its literals, and each
AND child contributes the singleton clauses of its literals.
\end{definition}

\begin{lemma}[DNF conversion preserves the value of a shallow OR gate]
\lean{LMN.depth2OrToDNF_eval}
\label{LMN.depth2OrToDNF_eval}
\leanok
\uses{BoolCircuit.Circuit, BoolCircuit.Circuit.depth, BoolCircuit.Circuit.eval, BoolCircuit.Lit, BoolCircuit.Lit.eval, BoolCircuit.Lit.eval_eq_toLiteral_eval, BoolCircuit.Lit.toLiteral, DNF.eval, LMN.depth2OrToDNF, LMN.depth_le_one_children_are_lits, LMN.depth_le_two_children_depth_le_one, Literal.eval, Term.eval}
\difficulty{5}
Let $cs$ be a finite list of Boolean circuits on $n$ variables, regarded as the children
of a top-level OR gate $\varphi$, and suppose $\varphi$ has depth at most $2$. Form a
DNF $D$ from $cs$ by turning each literal child into a singleton term, each AND child
into the single term consisting of all its literals, and each OR child into the
collection of singleton terms of its literals (in each case reading a circuit literal as
the corresponding DNF literal, keeping the same variable and polarity). Then $D$ and
$\varphi$ compute the same Boolean function: $D(x) = \varphi(x)$ for every assignment $x
\colon \{0,\dots,n-1\} \to \{\mathrm{true},\mathrm{false}\}$.
\end{lemma}

\begin{lemma}[Correctness of the depth-2 AND-to-CNF conversion]
\lean{LMN.depth2AndToCNF_eval}
\label{LMN.depth2AndToCNF_eval}
\leanok
\uses{BoolCircuit.Circuit, BoolCircuit.Circuit.depth, BoolCircuit.Circuit.eval, BoolCircuit.Lit, BoolCircuit.Lit.toLiteral, BoolCircuit.foldr_or_lits_eq_clause_eval, CNF.eval, CNF.evalClause, LMN.depth2AndToCNF, LMN.depth_le_two_children_depth_le_one, Literal.eval}
\difficulty{6}
Let $cs$ be a list of Boolean circuits on $n$ variables, viewed as the children of a
top-level AND gate, and suppose the AND-of-children circuit $\bigwedge cs$ has depth at
most $2$. Then the depth-$2$ AND-to-CNF conversion applied to $cs$ produces a CNF
formula computing the same Boolean function as $\bigwedge cs$: for every assignment $x :
\mathrm{Fin}\,n \to \mathrm{Bool}$, the value of the resulting CNF at $x$ equals the
value of $\bigwedge cs$ at $x$.
\end{lemma}

\begin{lemma}[Width bound for the depth-two OR-to-DNF conversion]
\lean{LMN.depth2OrToDNF_width_le}
\label{LMN.depth2OrToDNF_width_le}
\leanok
\uses{BoolCircuit.Circuit, BoolCircuit.Circuit.depth, BoolCircuit.Circuit.maxFanin, BoolCircuit.Lit.toLiteral, DNF.width, LMN.depth2OrToDNF}
\difficulty{5}
Let $cs$ be a list of Boolean circuits on $n$ variables, regarded as the children of a
top-level OR gate, and suppose that this OR gate has depth at most $2$. Form the DNF
obtained from $cs$ by converting each literal child into a singleton term, each AND
child into the term consisting of its literals, and each OR child into the singleton
terms of its literals. Then the width of this DNF (the maximum number of literals in any
of its terms) is at most the maximum fan-in of the OR gate.
\end{lemma}

\begin{lemma}[Width bound for the depth-two AND-to-CNF conversion]
\lean{LMN.depth2AndToCNF_width_le}
\label{LMN.depth2AndToCNF_width_le}
\leanok
\uses{BoolCircuit.Circuit, BoolCircuit.Circuit.depth, BoolCircuit.Circuit.maxFanin, BoolCircuit.Lit.toLiteral, CNF.width, LMN.depth2AndToCNF}
\difficulty{5}
Let $cs$ be a list of Boolean circuits on $n$ variables, regarded as the children of a
top AND gate, and suppose this AND gate has depth at most $2$. Then the CNF formula
obtained by converting these children — each literal child becoming a singleton clause,
each OR child the clause of its literals, and each AND child contributing the singleton
clauses of its literals — has width at most the maximum fanin of the AND gate.
\end{lemma}
